\section{Phương trình lượng giác cơ bản}

\subsection{Đề 1}
\Opensolutionfile{ans}[ans/ans-1-C1-B5.D1-CH]
\caulc
\begin{ex}%%[1D1N5-1]
	Nghiệm của phương trình lượng giác $\sin x=5$ là
	\choice
	{$x\in \mathbb{R}$}
	{$x=\pm \arcsin 5+k2\pi \left( k\in \mathbb{Z} \right)$}
	{$\left[ \begin{aligned}
		  & x=\arcsin 5+k2\pi  \\ 
		 & x=\pi -\arcsin 5+k2\pi  \\ 
		\end{aligned} \right.\left( k\in \mathbb{Z} \right)$}
	{\True $x\in \varnothing $}
	\loigiai{
		Phương trình $\sin x=m$ chỉ có nghiệm khi và chỉ khi $\left| m \right|\le 1$.\\
		Nên phương trình đã cho vô nghiệm.
	}
\end{ex}

\begin{ex}%[1D1H5-3]
	Nghiệm của phương trình $\cos x=-\dfrac{1}{2}$ là
	\choice
	{\True $x=\pm \dfrac{2\pi }{3}+k2\pi $}
	{$x=\pm \dfrac{\pi }{6}+k\pi $}
	{$x=\pm \dfrac{\pi }{3}+k2\pi $}
	{$x=\pm \dfrac{\pi }{6}+k2\pi $}
	\loigiai{
		Ta có $\cos x=-\dfrac{1}{2}\Leftrightarrow \cos x=\cos \left( -\dfrac{2\pi }{3}\right) \Leftrightarrow x=\pm \dfrac{2\pi }{3}+k2\pi.$\\
		Vậy phương trình có các nghiệm là $x=\pm \dfrac{2\pi }{3}+k2\pi \left( k\in \mathbb{Z} \right)$.
	}
\end{ex}

\begin{ex}%[1D1H5-3]
	Phương trình nào sau đây có nghiệm?
	\choice
	{$\cos x=-\dfrac{3}{2}$}
	{$\sin x=\sqrt{2}$}
	{\True $\tan x=3$}
	{${\cos }^{2}x-3=0$}
	\loigiai{
		Vì phương trình $\tan x=a$ luôn có nghiệm.
	}
\end{ex}

\begin{ex}%[1D1H5-3]
	Nghiệm của phương trình $\tan 2x+\sqrt{3}=0$ là
	\choice
	{$x=\dfrac{\pi }{6}+k\pi ,k\in \mathbb{Z}$}
	{$x=-\dfrac{\pi }{6}+k\pi ,k\in \mathbb{Z}$}
	{$x=\dfrac{\pi }{6}+\dfrac{k\pi }{2},k\in \mathbb{Z}$}
	{\True $x=-\dfrac{\pi }{6}+\dfrac{k\pi }{2},k\in \mathbb{Z}$}
	\loigiai{
		Ta có $\tan 2x+\sqrt{3}=0\Leftrightarrow \tan 2x=-\sqrt{3}\Leftrightarrow 2x=-\dfrac{\pi }{3}+k\pi $ $\Leftrightarrow x=-\dfrac{\pi }{6}+\dfrac{k\pi }{2},k\in \mathbb{Z}$.
	}
\end{ex}

\begin{ex}%[1D1H5-3]
	Nghiệm của phương trình $\sqrt{3}\cot x-3=0$ là	
	\choice
	{$x=\dfrac{\pi }{3}+k2\pi \left( k\in \mathbb{Z} \right)$}
	{$x=\dfrac{\pi }{3}+k\pi \left( k\in \mathbb{Z} \right)$}
	{\True $x=\dfrac{\pi }{6}+k\pi \left( k\in \mathbb{Z} \right)$}
	{$x=\dfrac{\pi }{6}+k2\pi \left( k\in \mathbb{Z} \right)$}
	\loigiai{
		Ta có $\sqrt{3}\cot x-3=0\Leftrightarrow \cot x=\sqrt{3}\Leftrightarrow x=\dfrac{\pi }{6}+k\pi \left( k\in \mathbb{Z} \right)$.
	}
\end{ex}

\begin{ex}%[1D1H5-3]
	Tìm tất cả các giá trị thực của tham số $m$để phương trình $\sin x=m$ có nghiệm thực.	
	\choice
	{$m\ge 0$}
	{\True $-1\le m\le 1$}
	{$-1<m<1$}
	{$m>0$}
	\loigiai{
		Do $-1\le \sin x\le 1,\text{ }\forall x\in \mathbb{R}$ nên phương trình $\sin x=m$ có nghiệm khi và chỉ khi $-1\le m\le 1$.
	}
\end{ex}

\begin{ex}%[1D1H5-3]
	Tập nghiệm của phương trình $3\cos\left( 3x-\dfrac{\pi }{3} \right)=0$ là
	\choice
	{$\left\{ \dfrac{\pi }{2}+k\pi ,k\in \mathbb{Z} \right\}$}
	{\True $\left\{ \dfrac{5\pi }{18}+\dfrac{k\pi }{3},k\in \mathbb{Z} \right\}$}
	{$\left\{ \dfrac{5\pi }{6}+k2\pi ,k\in \mathbb{Z} \right\}$}
	{$\left\{ \dfrac{5\pi }{18}+\dfrac{k2\pi }{3},k\in \mathbb{Z} \right\}$}
	\loigiai{
		Ta có $3\cos\left( 3x-\dfrac{\pi }{3} \right)=0\Leftrightarrow \cos\left( 3x-\dfrac{\pi }{3} \right)=0\Leftrightarrow 3x-\dfrac{\pi }{3}=\dfrac{\pi }{2}+k\pi \Leftrightarrow x=\dfrac{5\pi }{18}+\dfrac{k\pi }{3}$.
	}
\end{ex}

\begin{ex}%[1D1H5-3]
	Nghiệm của phương trình $\cos x=\dfrac{1}{2}$ là
	\choice
	{\True $\left[ \begin{aligned}
		  & x=\dfrac{\pi }{6}+k2\pi  \\ 
		 & x=-\dfrac{\pi }{6}+k2\pi . \\ 
		\end{aligned} \right.,k\in \mathbb{Z}$}
	{$\left[ \begin{aligned}
		  & x=\dfrac{\pi }{6}+k2\pi  \\ 
		 & x=\dfrac{5\pi }{6}+k2\pi  \\ 
		\end{aligned} \right.,k\in \mathbb{Z}$}
	{$\left[ \begin{aligned}
		  & x=\dfrac{\pi }{3}+k2\pi  \\ 
		 & x=\dfrac{2\pi }{3}+k2\pi  \\ 
		\end{aligned} \right.,k\in \mathbb{Z}$}
	{$\left[ \begin{aligned}
		  & x=\dfrac{\pi }{3}+k2\pi  \\ 
		 & x=-\dfrac{\pi }{3}+k2\pi  \\ 
		\end{aligned} \right.,k\in \mathbb{Z}$}
	\loigiai{
		Ta có $\cos x=\dfrac{1}{2}\Leftrightarrow \cos x=\cos \dfrac{\pi }{6}\Leftrightarrow \left[ \begin{aligned}
			  & x=\dfrac{\pi }{6}+k2\pi  \\ 
			 & x=-\dfrac{\pi }{6}+k2\pi  \\ 
			\end{aligned} \right.,k\in \mathbb{Z}$.
	}
\end{ex}

\begin{ex}%[1D1H5-3]
	Nghiệm của phương trình $\sqrt{3}+3\tan x=0$ là
	\choice
	{$x=\dfrac{\pi }{3}+k\pi $}
	{$x=\dfrac{\pi }{2}+k2\pi $}
	{\True $x=-\dfrac{\pi }{6}+k\pi $}
	{$x=\dfrac{\pi }{2}+k\pi $}
	\loigiai{
		Ta có $\sqrt{3}+3\tan x=0\Leftrightarrow \tan x=-\dfrac{\sqrt{3}}{3}\Leftrightarrow x=-\dfrac{\pi }{6}+k\pi \left( k\in \mathbb{Z} \right)$.
	}
\end{ex}

\begin{ex}%[1D1H5-4]
	Giải phương trình $\cot 2x=\cot 20^\circ $.	
	\choice
	{\True $x=10^\circ +k90^\circ ,k\in \mathbb{Z}$}
	{$x=10^\circ +k180^\circ ,k\in \mathbb{Z}$}
	{$x=20^\circ +k90^\circ ,k\in \mathbb{Z}$}
	{$x=20^\circ +k180^\circ ,k\in \mathbb{Z}$}
	\loigiai{
		Ta có $\cot 2x=\cot 20{}^\circ \Leftrightarrow 2x=20^\circ +k180^\circ \Leftrightarrow x=10{}^\circ +k90^\circ ,k\in \mathbb{Z}$.
	}
\end{ex}

\begin{ex}%[1D1H5-3]
	Phương trình $2\cot x-\sqrt{3}=0$ có nghiệm là
	\choice
	{\True $x=\text{arccot}\dfrac{\sqrt{3}}{2}+k\pi \left( k\in \mathbb{Z} \right)$}
	{$x=\dfrac{\pi }{6}+k\pi \left( k\in \mathbb{Z} \right)$}
	{$\left[ \begin{aligned}
		  & x=\dfrac{\pi }{6}+k2\pi  \\ 
		 & x=-\dfrac{\pi }{6}+k2\pi  \\ 
		\end{aligned} \right.\left( k\in \mathbb{Z} \right)$}
	{$x=\dfrac{\pi }{3}+k2\pi \left( k\in \mathbb{Z} \right)$}
	\loigiai{
		Ta có $2\cot x-\sqrt{3}=0\Leftrightarrow \cot x=\dfrac{\sqrt{3}}{2}\Leftrightarrow x=\text{arccot}\dfrac{\sqrt{3}}{2}+k\pi \left( k\in \mathbb{Z} \right)$.
	}
\end{ex}

\begin{ex}%[1D1H5-3]
	Tập nghiệm của phương trình $2\cos x+\sqrt{3}=0$ là
	\choice
	{\True $\left\{ \pm \dfrac{5\pi }{6}+k2\pi |k\in \mathbb{Z} \right\}$}
	{$\left\{ \pm \dfrac{\pi }{6}+k2\pi |k\in \mathbb{Z} \right\}$}
	{$\left\{ \pm \dfrac{\pi }{6}+k\pi |k\in \mathbb{Z} \right\}$}
	{$\left\{ \pm \dfrac{5\pi }{6}+k\pi |k\in \mathbb{Z} \right\}$}
	\loigiai{
		Ta có $2\cos x+\sqrt{3}=0\Leftrightarrow \cos x=-\dfrac{\sqrt{3}}{2}\Leftrightarrow x=\pm \dfrac{5\pi }{6}+k2\pi ,k\in \mathbb{Z}$.
	}
\end{ex}

\Closesolutionfile{ans}
\indapan{6}{ans/ans-1-C1-B5.D1-CH}

\cauds
\Opensolutionfile{ans}[ans/ans-1-C1-B5.D1-DS]
\begin{ex}%[1D1H5-3]
	Cho phương trình lượng giác $\sin 2x=-\dfrac{1}{2}\quad (*)$. Khi đó
	\choiceTF
	{Phương trình $(*)$ tương đương $\sin 2x=\sin \dfrac{\pi }{6}$}
	{Trong khoảng $\left( 0;\pi  \right)$ phương trình có $3$ nghiệm}
	{\True Tổng các nghiệm của phương trình trong khoảng $(0;\pi)$ bằng $\dfrac{3\pi }{2}$}
	{\True Trong khoảng $(0;\pi)$ phương trình có nghiệm lớn nhất bằng $\dfrac{11\pi }{12}$}
	\loigiai{
	\begin{itemchoice}
	\itemch \textbf{Sai}. Ta có $\sin 2x=-\dfrac{1}{2} \Leftrightarrow \sin 2x=\sin \left(-\dfrac{\pi}{6}\right)$.
	\itemch \textbf{Sai}. $\sin 2x=\sin \left(-\dfrac{\pi}{6}\right)\Leftrightarrow\hoac{& 2x=-\dfrac{\pi}{6}+k2\pi \\& 2x=\dfrac{7 \pi}{6}+k2 \pi}(k \in \mathbb{Z}) \Rightarrow\hoac{& x=-\dfrac{\pi}{12}+k \pi \\& x=\dfrac{7 \pi}{12}+k \pi}(k \in \mathbb{Z})$.
	\itemch \textbf{Đúng}. Với $0<x<\pi \Rightarrow \hoac{
		   & 0<\dfrac{-\pi }{12}+k\pi <\pi   \\
		   & 0<\dfrac{7\pi }{12}+k\pi <\pi}(k\in \mathbb{Z})\Leftrightarrow \hoac{
		   & k=1  \\
		   & k=0}\Leftrightarrow \hoac{
		   & x=\dfrac{11\pi }{12}  \\
		   & x=\dfrac{7\pi }{12}}$.\\
		   Tổng các nghiệm của phương trình trong khoảng $(0;\pi)$ bằng $\dfrac{3\pi }{2}$.
	\itemch \textbf{Đúng}.  Trong khoảng $(0;\pi)$ phương trình có nghiệm lớn nhất bằng $\dfrac{11\pi }{12}$.
	\end{itemchoice}
}
\end{ex}

\begin{ex}%[1D1H5-3]
	Cho phương trình lượng giác $\tan \left( 2x-15^{^\circ } \right)=1$ $(*)$. Khi đó:
	\choiceTF
	{\True Phương trình $(*)$ có nghiệm ${\tan \left(2 x-15^{\circ}\right)=1 \Leftrightarrow x=30^{\circ}+k 90^{\circ} (k \in \mathbb{Z})}$} 
	{Phương trình có nghiệm âm lớn nhất bằng $-30^\circ $}
	{Tổng các nghiệm của phương trình trong khoảng $\left(-180^\circ ;90^\circ  \right)$ bằng $180^\circ $}
	{Trong khoảng $\left(-180^\circ ;90^\circ  \right)$ phương trình có nghiệm lớn nhất bằng $60^\circ $}
	\loigiai{
		\begin{itemchoice}
		\itemch \textbf{Đúng}. Ta có ${\tan \left(2x-15^{\circ}\right)=1 \Leftrightarrow x=30^{\circ}+k 90^{\circ} (k \in \mathbb{Z})}$. 
		\itemch \textbf{Sai}. Với $k=-1\Rightarrow x=-60^{\circ}$. Phương trình có nghiệm âm lớn nhất bằng $-60^\circ $.
		\itemch \textbf{Sai}. Với $-180^{\circ}<x<90^{\circ} \Rightarrow-180^{\circ}<30^{\circ}+k 90^{\circ}<90^{\circ}(k \in \mathbb{Z}) \Rightarrow k=\{-2 ;-1 ; 0\}$.\\
		Vậy $\hoac{
		& x=-150^{\circ} \\
		& x=-60^{\circ} \\
		& x=30^{\circ}}$. Tổng các nghiệm của phương trình trong khoảng $\left(-180^\circ ;90^\circ  \right)$ bằng $-180^\circ $.
		\itemch \textbf{Sai}. Trong khoảng $\left( -180^\circ ;90^\circ  \right)$ phương trình có nghiệm lớn nhất bằng $30^\circ $.
		\end{itemchoice}
		}
\end{ex}

\begin{ex}%[1D1H5-3]
	Cho phương trình lượng giác $\cot 3x=-\dfrac{1}{\sqrt{3}}$ $(*)$. Khi đó
	\choiceTF
	{Phương trình $(*)$ tương đương $\cot 3x=\cot \left(\dfrac{-\pi }{6} \right)$}
	{Phương trình $(*)$ có nghiệm $x=\dfrac{\pi }{9}+k\dfrac{\pi }{3}(k\in \mathbb{Z})$}
	{\True Tổng các nghiệm của phương trình trong khoảng $\left( -\dfrac{\pi }{2};0 \right)$ bằng $\dfrac{-5\pi }{9}$}
	{\True Phương trình có nghiệm dương nhỏ nhất bằng $\dfrac{2\pi }{9}$}
	\loigiai{
		\begin{itemchoice}
		\itemch \textbf{Sai}. Ta có $\cot 3x=-\dfrac{1}{\sqrt{3}} \Leftrightarrow \cot 3x=\cot \left(\dfrac{-\pi}{3}\right)$.
		\itemch \textbf{Sai}. ${\cot 3 x=-\dfrac{1}{\sqrt{3}} \Leftrightarrow \cot 3x=\cot \left(\dfrac{-\pi}{3}\right) \Leftrightarrow 3x=\dfrac{-\pi}{3}+k \pi(k \in \mathbb{Z}) \Leftrightarrow x=\dfrac{-\pi}{9}+k \dfrac{\pi}{3}(k \in \mathbb{Z})}$.
		\itemch \textbf{Đúng}. Với $-\dfrac{\pi}{2}<x<0\Rightarrow-\dfrac{\pi}{2}<\dfrac{-\pi}{9}+k \dfrac{\pi}{3}<0(k \in \mathbb{Z})\Leftrightarrow \dfrac{-7}{6}<k<\dfrac{1}{3}$\\
		Vì $k\in \mathbb{Z}$ nên $k=\{-1 ; 0\} \Rightarrow\hoac{
				& x=\dfrac{-\pi}{9} \\
				& x=\dfrac{-4 \pi}{9}}$.\\
		Tổng các nghiệm của phương trình trong khoảng $\left( -\dfrac{\pi }{2};0 \right)$ bằng $\dfrac{-5\pi }{9}$.
		\itemch \textbf{Đúng}. Với $k=1\Rightarrow x=\dfrac{2\pi }{9}$. Phương trình có nghiệm dương nhỏ nhất bằng $\dfrac{2\pi }{9}$.
		\end{itemchoice}
}
\end{ex}

\begin{ex}%[1D1H5-3]
	Cho phương trình lượng giác ${3-\sqrt{3} \tan \left(2 x-\dfrac{\pi}{3}\right)=0}$ , khi đó:
	\choiceTF
	{Phương trình có nghiệm $x=\dfrac{\pi }{6}+\dfrac{k\pi }{2},k\in \mathbb{Z}$}
	{Phương trình có nghiệm âm lớn nhất bằng $-\dfrac{\pi }{3}$}
	{Khi $\dfrac{-\pi }{4}<x<\dfrac{2\pi }{3}$ thì phương trình có ba nghiệm}
	{\True Tổng các nghiệm của phương trình trong khoảng $\left( \dfrac{-\pi }{4};\dfrac{2\pi }{3} \right)$ bằng $\dfrac{\pi }{6}$}
	\loigiai{
		
		
		
\begin{itemchoice}
		\itemch \textbf{Sai}. Phương trình tương đương với ${\tan \left(2 x-\dfrac{\pi}{3}\right)=\sqrt{3} \Leftrightarrow x=\dfrac{\pi}{3}+\dfrac{k \pi}{2}, k \in \mathbb{Z}}$.
		\itemch \textbf{Sai}. Với $k=-1\Rightarrow x=-\dfrac{\pi }{6}$. Phương trình có nghiệm âm lớn nhất bằng $-\dfrac{\pi }{6}$.
		\itemch \textbf{Sai}. Vì $\dfrac{-\pi }{4}<x<\dfrac{2\pi }{3}\Leftrightarrow \dfrac{-\pi }{4}<\dfrac{\pi }{3}+\dfrac{k\pi }{2}<\dfrac{2\pi }{3}\Leftrightarrow \dfrac{-7\pi }{12}<\dfrac{k\pi }{2}<\dfrac{\pi }{3}\Leftrightarrow \dfrac{-7}{6}<k<\dfrac{2}{3}$
		Do ${k \in \mathbb{Z}}$ nên ${k \in\{-1 ; 0\}}$.
		\itemch \textbf{Đúng}. Với ${k=-1}$ thì ${x=\dfrac{-\pi}{6}}$, với ${k=0}$ thì $x=\dfrac{\pi}{3}$.\\
		Vậy tổng các nghiệm của phương trình trong khoảng $\left( \dfrac{-\pi }{4};\dfrac{2\pi }{3} \right)$ bằng $\dfrac{\pi }{6}$.
		\end{itemchoice}
}
\end{ex}
\Closesolutionfile{ans}
\indapan{2}{ans/ans-1-C1-B5.D1-DS}

\caukq
\Opensolutionfile{ans}[ans/ans-1-C1-B5.D1-KQ]
\begin{ex}%[1D1H5-3]
	Có tất cả bao nhiêu giá trị nguyên của $m$ để phương trình $\sqrt{3}\cos x+m-1=0$ có nghiệm?
	\shortans{3}
	\loigiai{	
		Ta có $\sqrt{3}\cos x+m-1=0\Leftrightarrow \cos x=\dfrac{1-m}{\sqrt{3}}$.\\
		Mặt khác $-1\le \cos x\le 1\Rightarrow -1\le \dfrac{1-m}{\sqrt{3}}\le 1\Leftrightarrow -\sqrt{3}\le 1-m\le \sqrt{3}\Leftrightarrow 1-\sqrt{3}\le m\le 1+\sqrt{3}$.\\
		Mà $m\in \mathbb{Z}$ nên $m\in \left\{ 1;2;3 \right\}$.	
	}
\end{ex}

\begin{ex}%[1D1H5-3]
	Tìm số điểm biểu diễn tất cả các nghiệm của phương trình $\sin \left( x+\dfrac{\pi }{3} \right)=\dfrac{1}{2}$ trên đường tròn lượng giác.
	\shortans{$2$}
	\loigiai{
		Ta có $\sin \left( x+\dfrac{\pi }{3} \right)=\dfrac{1}{2}\Leftrightarrow \left[ \begin{aligned}
		  & x+\dfrac{\pi }{3}=\dfrac{\pi }{6}+k2\pi  \\ 
		 & x+\dfrac{\pi }{3}=\dfrac{5\pi }{6}+k2\pi  \\ 
		\end{aligned} \right.\Leftrightarrow \left[ \begin{aligned}
		  & x=-\dfrac{\pi }{6}+k2\pi  \\ 
		 & x=\dfrac{\pi }{2}+k2\pi  \\ 
		\end{aligned} \right.\left( k\in \mathbb{Z} \right)$.\\
		Suy ra số điểm biểu diễn tất cả các nghiệm của phương trình đã cho trên đường tròn lượng giác là 2.
	}
\end{ex}

\begin{ex}%[1D1H5-3]
	Tổng tất cả các nghiệm của phương trình $3\cos x-1=0$trên đoạn $\left[ 0;4\pi  \right]$ là $a\pi$. Xác định $a$.
	\shortans{$8$}
	\loigiai{
		Ta có $3\cos x-1=0\Leftrightarrow \cos x=\dfrac{1}{3}\Leftrightarrow \left[ \begin{aligned}
		  & x=\alpha +k2\pi  \\ 
		 & x=-\alpha +k2\pi  \\ 
		\end{aligned} \right.$ (với $\alpha \in \left( 0;\dfrac{\pi }{2} \right)$, $k\in \mathbb{Z}$).\\
		Mà $x\in \left[ 0;4\pi  \right]$nên $x\in \left\{ \alpha ;-\alpha +2\pi ;\alpha +2\pi ;-\alpha +4\pi  \right\}$.\\
		Vậy tổng các nghiệm thỏa mãn đề bài là $\alpha +\left( -\alpha  \right)+2\pi +\alpha +2\pi +\left( -\alpha  \right)+4\pi =8\pi $. Vậy $a=8$.
}
\end{ex}

\begin{ex}%[1D1H5-3]
	Nghiệm dương nhỏ nhất của phương trình $\sin 5x+2{{\cos }^{2}}x=1$ có dạng $\dfrac{\pi a}{b}$ với $a$, $b$ là các số nguyên và nguyên tố cùng nhau. Tính tổng $S=a+b$.
	\shortans{$17$}
	\loigiai{
		Ta có $\sin 5x+2{{\cos }^{2}}x=1\Leftrightarrow \sin 5x=1-2{{\cos }^{2}}x\Leftrightarrow \sin 5x=-\cos 2x$ $\Leftrightarrow \sin 5x=\sin \left( 2x-\dfrac{\pi }{2} \right)$\\
		$\Leftrightarrow \left[ \begin{aligned}
		  & 5x=2x-\dfrac{\pi }{2}+k2\pi  \\ 
		 & 5x=\pi -\left( 2x-\dfrac{\pi }{2} \right)+k2\pi  \\ 
		\end{aligned} \right.\Leftrightarrow \left[ \begin{aligned}
		  & 3x=-\dfrac{\pi }{2}+k2\pi  \\ 
		 & 7x=\dfrac{3\pi }{2}+k2\pi  \\ 
		\end{aligned} \right.\Leftrightarrow \left[ \begin{aligned}
		  & x=-\dfrac{\pi }{6}+\dfrac{k2\pi }{3} \\ 
		 & x=\dfrac{3\pi }{14}+\dfrac{k2\pi }{7} \\ 
		\end{aligned} \right.\left( k\in \mathbb{Z} \right)$.\\
		Vì $x>0$ nên ta xét 2 trường hợp:\\
		Với $x=-\dfrac{\pi }{6}+\dfrac{k2\pi }{3}$ ta có $-\dfrac{\pi }{6}+\dfrac{k2\pi }{3}>0\Leftrightarrow 4k>1\Leftrightarrow k>\dfrac{1}{4}$ do $k\in \mathbb{Z}$ suy ra $k=1$ nên nghiệm dương nhỏ nhất trong trường hợp này là $x=\dfrac{\pi }{2}$.\\
		Với $x=\dfrac{3\pi }{14}+\dfrac{k2\pi }{7}$ ta có $\dfrac{3\pi }{14}+\dfrac{k2\pi }{7}>0$$\Leftrightarrow 4k>-3$ $\Leftrightarrow k>-\dfrac{3}{4}$ do $k\in \mathbb{Z}$ suy ra $k=0$ nên nghiệm dương nhỏ nhất trong trường hợp này là $x=\dfrac{3\pi }{14}$.\\
		Vậy nghiệm dương nhỏ nhất của phương trình là $x=\dfrac{3\pi }{14}$ suy ra $a=3$, $b=14$ $\Rightarrow S=17$.
}
\end{ex}

\begin{ex}%[1D1H5-3]
	Tìm số nghiệm của phương trình $\tan x=\tan \dfrac{3\pi }{8}$ trên $\left( \dfrac{\pi }{4};2\pi  \right)$.
	\shortans{$2$}
	\loigiai{
		Ta có $\tan x=\tan \dfrac{3\pi }{8}\Leftrightarrow x=\dfrac{3\pi }{8}+k\pi $, $k\in \mathbb{Z}$.\\
		Với $x\in \left( \dfrac{\pi }{4};2\pi  \right)$, ta có $\dfrac{\pi }{4}<\dfrac{3\pi }{8}+k\pi <2\pi \Leftrightarrow -\dfrac{1}{8}<k<\dfrac{13}{8}$ suy ra $k\in \left\{ 0;1 \right\}$.\\
		Vậy trên khoảng $\left( \dfrac{\pi }{4};2\pi  \right)$, phương trình đã cho có hai nghiệm.
}
\end{ex}

\begin{ex}%[1D1K5-6]
	Chiều cao $h$ (m) của một cabin trên vòng quay vào thời điểm $t$ giây sau khi bắt đầu chuyển động được cho bởi công thức $h(t)=30 + 20\sin \left(\dfrac{\pi}{25}t + \dfrac{\pi}{3}\right)$. Sau bao nhiêu giây thì cabin đạt độ cao $40$ m lần đầu tiên?
	\shortans{$12{,}5$}
	\loigiai{
		Để cabin đạt độ cao $40$ m thì 
		\begin{eqnarray*}
		& & h(t)=40 \Leftrightarrow 20\sin \left(\dfrac{\pi}{25}t + \dfrac{\pi}{3}\right) =10 \Leftrightarrow  \sin \left(\dfrac{\pi}{25}t + \dfrac{\pi}{3}\right) = \dfrac{1}{2}\\
		&\Leftrightarrow & \hoac{&  \dfrac{\pi}{25}t + \dfrac{\pi}{3} = \dfrac{\pi}{6} + k2\pi \\& \dfrac{\pi}{25}t + \dfrac{\pi}{3} = \dfrac{5\pi}{6} + k2\pi} \Leftrightarrow \hoac{& t = -\dfrac{25}{6}+50k \\ & t = \dfrac{25}{2} + 50k}, k \in \mathbb{Z}.
		\end{eqnarray*}
		Dễ thấy, giá trị $t>0$ nhỏ nhất thoả mãn là $t = 12{,}5$ (giây).
	}
\end{ex}
\Closesolutionfile{ans}
\indapan{6}{ans/ans-1-C1-B5.D1-KQ}
\subsection{Đề 2}
\Opensolutionfile{ans}[ans/ans-1-C1-B5.D2-CH]
\caulc
\begin{ex}%[1D1H5-3]
	Nghiệm của phương trình $\sin 2x=1$ là
	\choice
	{\True $x=\dfrac{\pi }{4}+k\pi $}
	{$x=\dfrac{\pi }{4}+k2\pi $}
	{$x=\dfrac{k\pi }{2}$}
	{$x=\dfrac{\pi }{2}+k2\pi $}
	\loigiai{
		Ta có $\sin 2x=1$$\Leftrightarrow 2x=\dfrac{\pi }{2}+k2\pi $$\Leftrightarrow x=\dfrac{\pi }{4}+k\pi $.
	}
\end{ex}

\begin{ex}%[1D1H5-3]
	Nghiệm của phương trình $\cos x=- \dfrac{1}{2}$ là
	\choice
	{\True $x=\pm \dfrac{2\pi }{3}+k2\pi $}
	{$x=\pm \dfrac{\pi }{6}+k\pi $}
	{$x=\pm \dfrac{\pi }{3}+k2\pi $}
	{$x=\pm \dfrac{\pi }{6}+k2\pi $}
	\loigiai{
		Ta có $\cos x=- \dfrac{1}{2}\Leftrightarrow \cos x=\cos \dfrac{2\pi }{3}\Leftrightarrow x=\pm \dfrac{2\pi }{3}+k2\pi ,k\in \mathbb{Z}$.
	}
\end{ex}

\begin{ex}%[1D1H5-3]
	Phương trình $\cos x=-\dfrac{\sqrt{3}}{2}$ có tập nghiệm là
	\choice
	{$\left\{ x=\pm \dfrac{\pi }{6}+k\pi ; k\in \mathbb{Z} \right\}$}
	{\True $\left\{ x=\pm \dfrac{5\pi }{6}+k2\pi ; k\in \mathbb{Z} \right\}$}
	{$\left\{ x=\pm \dfrac{\pi }{3}+k\pi ; k\in \mathbb{Z} \right\}$}
	{$\left\{ x=\pm \dfrac{\pi }{3}+k2\pi ; k\in \mathbb{Z} \right\}$}
	\loigiai{
		Ta có $\cos x=-\dfrac{\sqrt{3}}{2}$$\Leftrightarrow \cos x=\cos \dfrac{5\pi }{6}$$\Leftrightarrow x=\pm \dfrac{5\pi }{6}+k2\pi ,k\in \mathbb{Z}$.
	}
\end{ex}

\begin{ex}%[1D1H5-3]
	Giải phương trình $\sqrt{\text{3}}\tan 2x-3=0$.
	\choice
	{$x=\dfrac{\pi }{3}+k\dfrac{\pi }{2}\left( k\in \mathbb{Z} \right)$}
	{$x=\dfrac{\pi }{3}+k\pi \left( k\in \mathbb{Z} \right)$}
	{\True $x=\dfrac{\pi }{6}+k\dfrac{\pi }{2}\left( k\in \mathbb{Z} \right)$}
	{$x=\dfrac{\pi }{6}+k\pi \left( k\in \mathbb{Z} \right)$}
	\loigiai{
		Ta có $\sqrt{\text{3}}\tan 2x-3=0\Leftrightarrow \tan 2x=\sqrt{3}$$\Leftrightarrow 2x=\dfrac{\pi }{3}+k\pi $$\Leftrightarrow x=\dfrac{\pi }{6}+k\dfrac{\pi }{2}\left( k\in \mathbb{Z} \right)$.
	}
\end{ex}

\begin{ex}%[1D1H5-3]
	Tập nghiệm của phương trình $\tan x+1=0$ là
	\choice
	{$S=\left\{ -\dfrac{\pi }{4}+k2\pi ,k\in \mathbb{Z} \right\}$}
	{$S=\left\{ \dfrac{\pi }{4}+k\pi ,k\in \mathbb{Z} \right\}$}
	{\True $S=\left\{ -\dfrac{\pi }{4}+k\pi ,k\in \mathbb{Z} \right\}$}
	{$S=\left\{ \dfrac{\pi }{4}+k2\pi ,k\in \mathbb{Z} \right\}$}
	\loigiai{
		Ta có $\tan x+1=0$$\Leftrightarrow \tan x=-1$ $\Leftrightarrow x=-\dfrac{\pi }{4}+k\pi ,k\in \mathbb{Z}$.
	}
\end{ex}

\begin{ex}%[1D1H5-3]
	Phương trình $\cot \left( x+45^\circ  \right)=\dfrac{\sqrt{3}}{3}$ có nghiệm là (với $k\in \mathbb{Z}$)
	\choice
	{\True $15^\circ +k180{}^\circ $}
	{$30^\circ +k180{}^\circ $}
	{$45^\circ +k180^\circ $}
	{$60^\circ +k180^\circ $}
	\loigiai{
		Phương trình $\cot \left( x+45^\circ  \right)=\dfrac{\sqrt{3}}{3}$ $\Leftrightarrow \cot \left( x+45^\circ  \right)=\cot 60{}^\circ $ $\Leftrightarrow x+45{}^\circ =60^\circ +k180^\circ $
		$\Leftrightarrow x=15{}^\circ +k180^\circ $ ( với $k\in \mathbb{Z}$).
	}
\end{ex}

\begin{ex}%[1D1H5-3]
	Phương trình $\cos x-m=0$ vô nghiệm khi giá trị tham số $m$thỏa mãn.
	\choice
	{\True $\left[ \begin{aligned}
		  & m<-1 \\ 
		 & m>1 \\ 
		\end{aligned} \right.$}
	{$-1\le m\le 1$}
	{$m>1$}
	{$m<-1$}
	\loigiai{
		Ta có $\cos x-m=0\Leftrightarrow \cos x=m$.
			Phương trình vô nghiệm khi và chỉ khi $m\notin \left[ -1;1 \right]$$\Leftrightarrow \left[ \begin{aligned}
			  & m<-1 \\ 
			 & m>1 \\ 
			\end{aligned} \right..$
	}
\end{ex}

\begin{ex}%[1D1H5-3]
	Phương trình $\cos x=-\dfrac{\sqrt{3}}{2}$ có tập nghiệm là
	\choice
	{$\left\{ x=\pm \dfrac{\pi }{6}+k\pi ; k\in \mathbb{Z} \right\}$}
	{\True $\left\{ x=\pm \dfrac{5\pi }{6}+k2\pi ; k\in \mathbb{Z} \right\}$}
	{$\left\{ x=\pm \dfrac{\pi }{3}+k\pi ; k\in \mathbb{Z} \right\}$}
	{$\left\{ x=\pm \dfrac{\pi }{3}+k2\pi ; k\in \mathbb{Z} \right\}$}
	\loigiai{
		Ta có $\cos x=-\dfrac{\sqrt{3}}{2}$$\Leftrightarrow \cos x=\cos \dfrac{5\pi }{6}$$\Leftrightarrow x=\pm \dfrac{5\pi }{6}+k2\pi ,k\in \mathbb{Z}$.
	}
\end{ex}

\begin{ex}%[1D1H5-4]
	Giải phương trình $\tan (2x)=\tan 80^\circ $. Kết quả thu được là
	\choice
	{$x=80^\circ +k180^\circ $}
	{\True $x=40^\circ +k90^\circ $}
	{$x=40^\circ +k45^\circ $}
	{$x=40^\circ +k180^\circ $}
	\loigiai{
		Ta có $\tan \left( 2x \right)=\tan 80^\circ $$\Leftrightarrow 2x=80^\circ +k180^\circ $$\Leftrightarrow x=40^\circ +k90^\circ ,k\in \mathbb{Z}$.
	}
\end{ex}

\begin{ex}%[1D1H5-3]
	Tập nghiệm của phương trình $\cot 2x=0$ là
	\choice
	{$\left\{ \dfrac{\pi }{4}+k2\pi |k\in \mathbb{Z} \right\}$}
	{$\left\{ \dfrac{\pi }{4}+k\pi |k\in \mathbb{Z} \right\}$}
	{\True $\left\{ \dfrac{\pi }{4}+k\dfrac{\pi }{2}|k\in \mathbb{Z} \right\}$}
	{$\left\{ \dfrac{\pi }{4}+k4\pi |k\in \mathbb{Z} \right\}$}
	\loigiai{
		Ta có $\cot 2x=0\Leftrightarrow 2x=\dfrac{\pi }{2}+k\pi \Leftrightarrow x=\dfrac{\pi }{4}+k\dfrac{\pi }{2}(k\in \mathbb{Z}).$
	}
\end{ex}

\begin{ex}%[1D1H5-3]
	Nghiệm của phương trình $\cos \left( x+\dfrac{\pi }{6} \right)=-\dfrac{1}{2}$ là
	\choice
	{$\left[ \begin{aligned}
		  & x=\dfrac{\pi }{2}+k2\pi  \\ 
		 & x=\dfrac{\pi }{3}+k2\pi  \\ 
		\end{aligned} \right.\left( k\in \mathbb{Z} \right)$}
	{\True $\left[ \begin{aligned}
		  & x=\dfrac{\pi }{2}+k2\pi  \\ 
		 & x=-\dfrac{5\pi }{6}+k2\pi  \\ 
		\end{aligned} \right.\left( k\in \mathbb{Z} \right)$}
	{$\left[ \begin{aligned}
		  & x=\dfrac{\pi }{2}+k2\pi  \\ 
		 & x=\dfrac{\pi }{6}+k2\pi  \\ 
		\end{aligned} \right.\left( k\in \mathbb{Z} \right)$}
	{$\left[ \begin{aligned}
		  & x=\dfrac{\pi }{6}+k2\pi  \\ 
		 & x=-\dfrac{5\pi }{6}+k2\pi  \\ 
		\end{aligned} \right.\left( k\in \mathbb{Z} \right)$}
	\loigiai{
		Ta có $\cos \left( x+\dfrac{\pi }{6} \right)=-\dfrac{1}{2}\Leftrightarrow \left[ \begin{aligned}
		  & x+\dfrac{\pi }{6}=\dfrac{2\pi }{3}+k2\pi  \\ 
		 & x+\dfrac{\pi }{6}=-\dfrac{2\pi }{3}+k2\pi  \\ 
		\end{aligned} \right.\Leftrightarrow \left[ \begin{aligned}
		  & x=\dfrac{\pi }{2}+k2\pi  \\ 
		 & x=-\dfrac{5\pi }{6}+k2\pi  \\ 
		\end{aligned} \right.\left( k\in \mathbb{Z} \right)$.
	}
\end{ex}

\begin{ex}%[1D1H5-3]
	Tìm tập nghiệm $S$ của phương trình $\cos 2x=-\dfrac{\sqrt{2}}{2}$.
	\choice
	{$S=\left\{ -\dfrac{3\pi }{8}+k2\pi ;\dfrac{3\pi }{8}+k2\pi ,k\in \mathbb{Z} \right\}$}
	{\True $S=\left\{ -\dfrac{3\pi }{8}+k\pi ;\dfrac{3\pi }{8}+k\pi ,k\in \mathbb{Z} \right\}$}
	{$S=\left\{ \dfrac{3\pi }{8}+k\pi ;\dfrac{\pi }{8}+k\pi ,k\in \mathbb{Z} \right\}$}
	{$S=\left\{ \dfrac{3\pi }{8}+k2\pi ;\dfrac{\pi }{8}+k2\pi ,k\in \mathbb{Z} \right\}$}
	\loigiai{
		Ta có $\cos 2x=-\dfrac{\sqrt{2}}{2}\Leftrightarrow \left[ \begin{aligned}
			  & 2x=\dfrac{3\pi }{4}+k2\pi  \\ 
			 & 2x=-\dfrac{3\pi }{4}+k2\pi  \\ 
			\end{aligned} \right.\Leftrightarrow \left[ \begin{aligned}
			  & x=\dfrac{3\pi }{8}+k\pi  \\ 
			 & x=-\dfrac{3\pi }{8}+k\pi  \\ 
			\end{aligned} \right.\left( k\in \mathbb{Z} \right).$\\
		Vậy tập nghiệm của phương trình trên là $S=\left\{ -\dfrac{3\pi }{8}+k\pi ;\dfrac{3\pi }{8}+k\pi ,k\in \mathbb{Z} \right\}.$
	}
\end{ex}

\Closesolutionfile{ans}
\indapan{6}{ans/ans-1-C1-B5.D2-CH}

\cauds
\Opensolutionfile{ans}[ans/ans-1-C1-B5.D2-DS]
\begin{ex}%[1D1H5-3]
	Cho phương trình lượng giác $2 \cos x=\sqrt{3}$, khi đó
	\choiceTF
	{Phương trình có nghiệm $x=\pm \dfrac{\pi }{3}+k2\pi (k\in \mathbb{Z})$}
	{Trong đoạn ${x \in\left[0 ; \dfrac{5 \pi}{2}\right]}$ phương trình có 4 nghiệm}
	{\True Tổng các nghiệm của phương trình trong đoạn ${x \in\left[0 ; \dfrac{5 \pi}{2}\right]}$ bằng $\dfrac{25\pi }{6}$}
	{\True Trong đoạn ${x \in\left[0 ; \dfrac{5 \pi}{2}\right]}$ phương trình có nghiệm lớn nhất bằng $\dfrac{13\pi }{6}$}
	\loigiai{
		\begin{itemchoice}
		\itemch \textbf{Sai}. Ta có ${2 \cos x=\sqrt{3} \Leftrightarrow \cos x=\dfrac{\sqrt{3}}{2} \Leftrightarrow x= \pm \dfrac{\pi}{6}+k 2 \pi(k \in \mathbb{Z})}$.
		\itemch \textbf{Sai}. Vì ${x \in\left[0 ; \dfrac{5 \pi}{2}\right]}$ nên ${x \in\left\{\dfrac{\pi}{6} ; \dfrac{11 \pi}{6} ; \dfrac{13 \pi}{6}\right\}}$.
		\itemch \textbf{Đúng}. Tổng các nghiệm của phương trình trong đoạn ${x \in\left[0 ; \dfrac{5 \pi}{2}\right]}$ bằng $\dfrac{25\pi }{6}$.
		\itemch \textbf{Đúng}. Trong đoạn ${x \in\left[0 ; \dfrac{5 \pi}{2}\right]}$ phương trình có nghiệm lớn nhất bằng $\dfrac{13\pi }{6}$.
		\end{itemchoice}
}
\end{ex}

\begin{ex}%[1D1H5-3]
	Cho phương trình ${\sin \left(2 x-\dfrac{\pi}{4}\right)=\sin \left(x+\dfrac{3 \pi}{4}\right)}\quad (*)$. Khi đó
	\choiceTF
	{\True Phương trình có nghiệm $\hoac{
		   & x=\pi +k2\pi   \\
		   & x=\dfrac{\pi }{6}+k\dfrac{2\pi }{3}}(k\in \mathbb{Z})$}
	{\True Trong khoảng $(0;\pi )$ phương trình có 2 nghiệm}
	{Tổng các nghiệm của phương trình trong khoảng $(0;\pi )$ bằng $\dfrac{7\pi }{6}$}
	{\True Trong khoảng $(0;\pi )$ phương trình có nghiệm lớn nhất bằng $x=\dfrac{5 \pi}{6}$}
	\loigiai{
		\begin{itemchoice}
		\itemch \textbf{Sai}. Ta có $\sin \left(2 x-\dfrac{\pi}{4}\right)=\sin \left(x+\dfrac{3 \pi}{4}\right) \Leftrightarrow\hoac{& 2 x-\dfrac{\pi}{4}=x+\dfrac{3 \pi}{4}+k 2 \pi \\ & 2  x-\dfrac{\pi}{4}=\dfrac{\pi}{4}-x+k 2 \pi}\Leftrightarrow\hoac{
		&x=\pi+k 2 \pi \\
		&x=\dfrac{\pi}{6}+k \dfrac{2 \pi}{3}}(k \in \mathbb{Z})$.
		\itemch \textbf{Sai}. Vì $x \in(0 ; \pi)$ nên $x \in\left\{\dfrac{\pi}{6} ; \dfrac{5 \pi}{6}\right\}$.\\
		Vậy phương trình có hai nghiệm thuộc khoảng ${(0 ; \pi)}$ là $x=\dfrac{\pi}{6}; x=\dfrac{5 \pi}{6}$.
		\itemch \textbf{Đúng}. Tổng các nghiệm của phương trình trong khoảng $(0;\pi )$ bằng $\dfrac{7\pi }{6}$.
		\itemch \textbf{Đúng}. Trong khoảng $(0;\pi )$ phương trình có nghiệm lớn nhất bằng $x=\dfrac{5 \pi}{6}$.
		\end{itemchoice}
}
\end{ex}

\begin{ex}%[1D1H5-3]
	Cho phương trình lượng giác ${\sin \left(3 x+\dfrac{\pi}{3}\right)=-\dfrac{\sqrt{3}}{2}}$.
	\choiceTF
	{Phương trình có nghiệm $\hoac{
	   & x=-\dfrac{\pi }{9}+k\dfrac{2\pi }{3}  \\
	   & x=\dfrac{\pi }{3}+k\dfrac{2\pi }{3} }(k\in \mathbb{Z})$}
	{\True Phương trình có nghiệm âm lớn nhất bằng $-\dfrac{2\pi }{9}$}
	{Trên khoảng $\left( 0;\dfrac{\pi }{2} \right)$ phương trình đã cho có 3 nghiệm}
	{\True Tổng các nghiệm của phương trình trong khoảng $\left( 0;\dfrac{\pi }{2} \right)$ bằng $\dfrac{7\pi }{9}$}
	\loigiai{
		\begin{itemchoice}
		\itemch \textbf{Sai}. Ta có 
		\begin{eqnarray*}
		& & \sin \left(3x+\dfrac{\pi}{3}\right)=-\dfrac{\sqrt{3}}{2}\\ &\Leftrightarrow&\left[\begin{array}{l}3 x+\dfrac{\pi}{3}=-\dfrac{\pi}{3}+k 2 \pi \\ 3 x+\dfrac{\pi}{3}=\dfrac{4 \pi}{3}+k 2 \pi\end{array}\right.\\
		&\Leftrightarrow & \hoac{& 3x=-\dfrac{2 \pi}{3}+k 2 \pi \\& 3x=\pi+k 2 \pi}\\ 
		&\Leftrightarrow&\hoac{& x=-\dfrac{2 \pi}{9}+k \dfrac{2 \pi}{3} \\& x=\dfrac{\pi}{3}+k \dfrac{2 \pi}{3}(k \in \mathbb{Z})}.
		\end{eqnarray*}
		\itemch \textbf{Sai}. Vì ${x \in\left(0 ; \dfrac{\pi}{2}\right)}$ nên ${x=\dfrac{\pi}{3}, x=\dfrac{4 \pi}{9}}$.
		\itemch \textbf{Đúng}. Vậy phương trình đã cho có hai nghiệm thuộc khoảng ${\left(0 ; \dfrac{\pi}{2}\right)}$..
		\itemch \textbf{Đúng}. Tổng các nghiệm của phương trình trong khoảng $\left( 0;\dfrac{\pi }{2} \right)$ bằng $\dfrac{7\pi }{9}$.
		\end{itemchoice}
}
\end{ex}

\begin{ex}%[1D1H5-3]
	Cho phương trình lượng giác ${2 \sin \left(x-\dfrac{\pi}{12}\right)+\sqrt{3}=0}$, khi đó
	\choiceTF
	{Phương trình tương đương $\sin \left( x-\dfrac{\pi }{12} \right)=\sin \left( \dfrac{\pi }{3} \right)$}
	{Phương trình có nghiệm là: $x=\dfrac{\pi }{4}+k2\pi ;x=\dfrac{7\pi }{12}+k2\pi (k\in \mathbb{Z})$}
	{\True Phương trình có nghiệm âm lớn nhất bằng $-\dfrac{\pi }{4}$}
	{\True Số nghiệm của phương trình trong khoảng $\left( -\pi ;\pi  \right)$ là hai nghiệm}
	\loigiai{
		\begin{itemchoice}
		\itemch \textbf{Sai}. Ta có 
		\begin{eqnarray*}
		& & 2 \sin \left(x-\dfrac{\pi}{12}\right)+\sqrt{3}=0\\ &\Leftrightarrow& \sin \left(x-\dfrac{\pi}{12}\right)=-\dfrac{\sqrt{3}}{2}\\ &\Leftrightarrow& \sin \left(x-\dfrac{\pi}{12}\right)=\sin \left(-\dfrac{\pi}{3}\right)\\
		&\Leftrightarrow & \hoac{
					& x-\dfrac{\pi}{12} = -\dfrac{\pi}{3}+k2\pi \\
					& x-\dfrac{\pi}{12} = \pi -(-\dfrac{\pi}{3})+k2\pi}\\
					&\Leftrightarrow& \hoac{
					& x=-\dfrac{\pi}{4}+k 2 \pi \\
					& x=\dfrac{17 \pi}{12}+k 2 \pi}(k \in \mathbb{Z})
		\end{eqnarray*}
		\itemch \textbf{Sai}. Vậy phương trình có nghiệm là ${x=-\dfrac{\pi}{4}+k 2 \pi ; x=\dfrac{17 \pi}{12}+k 2 \pi(k \in \mathbb{Z})}$..
		\itemch \textbf{Đúng}. Phương trình có nghiệm âm lớn nhất bằng $-\dfrac{\pi }{4}$.
		\itemch \textbf{Đúng}. Số nghiệm của phương trình trong khoảng $\left( -\pi ;\pi  \right)$ là hai nghiệm.
		\end{itemchoice}
}
\end{ex}


\Closesolutionfile{ans}
\indapan{2}{ans/ans-1-C1-B5.D2-DS}

\caukq
\Opensolutionfile{ans}[ans/ans-1-C1-B5.D2-KQ]
\begin{ex}%[1D1H5-3]
	Có bao nhiêu số nguyên $m$ để phương trình $3\cos \left( x+\dfrac{\pi }{6} \right)-m+5=0$ có nghiệm?
	\shortans{$7$}
	\loigiai{
		Ta có $3\cos \left( x+\dfrac{\pi }{6} \right)-m+5=0\Leftrightarrow \cos \left( x+\dfrac{\pi }{6} \right)=\dfrac{m-5}{3}$.\\
		Điều kiện để phương trình có nghiệm: $-1\le \dfrac{m-5}{3}\le 1\Leftrightarrow 2\le m\le 8$.\\
		Do $m$ nguyên nên $m=\left\{ 2;3;4;5;6;7;8 \right\}$. Vậy có $7$ số nguyên $m$.
	}
\end{ex}

\begin{ex}%[1D1H5-3]
	Số nghiệm của phương trình $\sin 2x=0$ thỏa mãn $0<x<2\pi $.
	\shortans{$3$}
	\loigiai{
		Ta có $\sin 2x=0$$\Leftrightarrow 2x=k\pi \Leftrightarrow x=k\dfrac{\pi }{2}\,\,\left( k\in \mathbb{Z} \right)$.\\
		Ta có $0<x<2\pi \Leftrightarrow 0<k\dfrac{\pi }{2}<2\pi $$\Leftrightarrow 0<k<4$.\\
		Mà $k$ nguyên nên $k\in \left\{ 1\,;\,2\,;\,3 \right\}\Rightarrow $ có $3$ giá trị của $k$ thỏa mãn .\\
		Vậy phương trình đã cho có $3$ nghiệm thuộc $(0;2\pi)$.
	}
\end{ex}

\begin{ex}%[1D1H5-3]
	Phương trình $\sin \left( x+\dfrac{\pi }{4} \right)+\cos x=0$ có tập nghiệm được biểu diễn bởi bao nhiêu điểm trên đường tròn lượng giác?	
	\shortans{$2$}
	\loigiai{
		Ta có 
		\begin{eqnarray*}
		& & \sin \left( x+\dfrac{\pi }{4} \right)+\cos x=0\\
		&\Leftrightarrow & \sin \left( x+\dfrac{\pi }{4} \right)=\sin \left( x-\dfrac{\pi }{2} \right)\\
		&\Leftrightarrow & x+\dfrac{\pi }{4}=\pi -x+\dfrac{\pi }{2}+\,k2\pi \\
		&\Leftrightarrow & 2x=\dfrac{5\pi }{4}+k2\pi\\
		&\Leftrightarrow & x=\dfrac{5\pi }{8}+k\pi.
		\end{eqnarray*}
		Ta có $0\leq x<2\pi \Leftrightarrow 0\leq \dfrac{5\pi }{8}+k\pi<2\pi\Leftrightarrow -\dfrac{5}{8}\leq k \dfrac{11}{8}$ mà $k\in \mathbb{Z}\Rightarrow k\in\{0;1\}$.\\
		Vậy phương trình đã cho có tập nghiệm biểu diễn được hai điểm trên đường tròn lương giác.
	}
\end{ex}

\begin{ex}%[1D1H5-3]
	Số nghiệm thực của phương trình $2\sin x+1=0$ trên đoạn $\left[ -\dfrac{3\pi }{2};\,10\pi  \right]$.
	\shortans{$12$}
	\loigiai{
		Phương trình tương đương $\sin x=-\dfrac{1}{2}$$\Leftrightarrow \left[ \begin{aligned}
			  & x=\dfrac{-\pi }{6}+k2\pi  \\ 
			 & x=\dfrac{7\pi }{6}+k2\pi  \\ 
			\end{aligned} \right.$, ($k\in \mathbb{Z}$).\\
			+ Với $x=-\dfrac{\pi }{6}+k2\pi $, $k\in \mathbb{Z}$ ta có $-\dfrac{3\pi }{2}\le -\dfrac{\pi }{6}+k2\pi \le 10\pi \Leftrightarrow \dfrac{-2}{3}\le k\le \dfrac{61}{12}$, $k\in \mathbb{Z}$
			$\Rightarrow 0\le k\le 5$.\\
			Do đó phương trình có $6$ nghiệm.\\
			+ Với $x=\dfrac{7\pi }{6}+k2\pi $, $k\in \mathbb{Z}$ ta có $-\dfrac{3\pi }{2}\le \dfrac{7\pi }{6}+k2\pi \le 10\pi \Leftrightarrow \dfrac{-4}{3}\le k\le \dfrac{53}{12}$, $k\in \mathbb{Z}$
			$\Rightarrow -1\le k\le 4$.\\
			Do đó, phương trình có $6$ nghiệm.\\
			+ Rõ ràng các nghiệm này khác nhau từng đôi một, vì nếu
			$-\dfrac{\pi }{6}+k2\pi =\dfrac{7\pi }{6}+{k}'2\pi \Leftrightarrow k-{k}'=\dfrac{2}{3}$ (vô lí, do $k$, ${k}'\in \mathbb{Z}$).\\
			Vậy phương trình có $12$ nghiệm trên đoạn $\left[ -\dfrac{3\pi }{2};\,10\pi  \right]$.
	}
\end{ex}

\begin{ex}%[1D1H5-3]
	Gọi $n$là số nghiệm của phương trình $\sin \left( 2x+30^\circ \right)=\dfrac{\sqrt{3}}{2}$ trên khoảng $\left( -180^\circ ;180^\circ\right)$. Tìm $n$.
	\shortans{$4$}
	\loigiai{
		Ta có $\sin \left( 2x+30^\circ \right)=\dfrac{\sqrt{3}}{2}\Leftrightarrow \left[ \begin{aligned}
		  & 2x+30^\circ=60^\circ+k360^\circ \\ 
		 & 2x+30^\circ=120^\circ+k360^\circ \\ 
		\end{aligned} \right.\Leftrightarrow \left[ \begin{aligned}
		  & x=15^\circ+k180^\circ \\ 
		 & x=45^\circ+k180^\circ \\ 
		\end{aligned} \right.\quad (k\in \mathbb{Z}).$\\
		Do $x\in \left( -180^\circ;180^\circ\right)$ nên $x\in \left\{ 15^\circ;-165^\circ;\,45^\circ;-135^\circ \right\}$. Vậy $n=4$.	
}
\end{ex}

\begin{ex}%[1D1K5-6]
	Một quả bóng được ném xiên một góc $\alpha$ $(0^\circ \le \alpha \le 90^\circ)$ từ mặt đất với tốc độ $v_0$ (m/s). Khoảng cách theo phương ngang từ vị trí ban đầu của quá bóng đến vị trí bóng chạm đất được tính bởi công thức $d = \dfrac{v_0^2 \sin 2\alpha}{10}$. Tính khoảng cách $d$ khi bóng được ném đi với tốc độ ban đầu $10$ m/s và góc ném là $30^\circ$ so với phương ngang (làm tròn đến hàng phần trăm).
	\shortans{$8{,}66$}
	\loigiai{
		Khi bóng được ném đi với tốc độ ban đầu $10$ m/s và góc ném là $30^\circ$ so với phương ngang thì 
		$$d = \dfrac{10^2 \sin 60^\circ}{10} = 5\sqrt{3} \approx 8{,}66.$$
		
	}
\end{ex}

\Closesolutionfile{ans}
\indapan{6}{ans/ans-1-C1-B5.D2-KQ}
